\chapter{Footprinting \& Network Scanning}
\label{ch:footprinting}

Once passive recon is done, active scanning starts. This phase maps the live
hosts, open ports, and running services. Nmap is the primary tool. Everything
feeds into enumeration and eventually exploitation.

% ─────────────────────────────────────────────────────────────────────────────
\section{Network Mapping --- Finding Live Hosts}
\label{sec:network_mapping}

Before scanning ports, I confirm which hosts are actually alive on the network.
Scanning ports on dead hosts wastes time.

\begin{termbox}
# Find your own IP to understand the network range
ip a s
ifconfig

# Ping sweep to find live hosts in a range
nmap -sn 192.168.1.0/24

# Scan a list of IPs from a file
nmap -sn -iL targets.txt
\end{termbox}

\screenshot{assets/images/Screenshot_2025-07-17_at_4.42.11_PM.png}{Nmap ping sweep identifying live hosts on the network}

% ─────────────────────────────────────────────────────────────────────────────
\section{Host Discovery with Nmap}
\label{sec:host_discovery}

\subsection{Basic Port Scan}

\begin{termbox}
nmap target_ip                    # Default scan (top 1000 ports)
nmap -p- target_ip                # All 65535 ports (slow but thorough)
nmap -p 22,80,443,445 target_ip   # Specific ports only
\end{termbox}

\subsection{Service and Version Detection}

\begin{whybox}[Why -sV matters]
Knowing a port is open is not enough. Version detection tells me the exact software
running, which I can cross-reference with exploit databases. Port 445 open is
interesting --- port 445 running SMBv1 on Windows 7 is a confirmed EternalBlue target.
\end{whybox}

\begin{termbox}
nmap -sV target_ip               # Version detection
nmap -sV -sC target_ip           # Version + default NSE scripts
nmap -sV -sC -p- target_ip       # All ports, version, scripts
nmap -sV -sC -O target_ip        # Add OS detection
\end{termbox}

\screenshot{assets/images/Screenshot_2025-07-19_at_3.38.46_PM.png}{Nmap -sV -sC output showing services and versions}

\screenshot{assets/images/Screenshot_2025-07-19_at_3.43.55_PM.png}{Nmap OS detection and script results}

\subsection{Stealth Scan}

\begin{termbox}
nmap -sS target_ip               # SYN (stealth) scan --- does not complete TCP handshake
nmap -f target_ip                # Packet fragmentation (may evade basic filters)
\end{termbox}

\begin{notebox}
\ic{-sS} sends SYN packets and waits for SYN-ACK without completing the handshake.
Faster than a full connect scan and less likely to be logged by the application layer,
though still visible at the network level.
\end{notebox}

\subsection{Metasploit Integration --- db\_nmap}

Nmap can be run from inside Metasploit. Results are automatically saved to the
workspace database, making them queryable with \ic{hosts}, \ic{services}, and
\ic{vulns}.



\begin{termbox}
# Inside msfconsole
service postgresql start && msfconsole
workspace -a project_name         # Create isolated workspace
setg RHOSTS target_ip             # Set target globally (all modules inherit this)
setg RHOST target_ip

db_nmap -sS -sV -o target_ip     # Run nmap, save to DB
hosts                             # View discovered hosts
services                          # View discovered services
\end{termbox}

\begin{notebox}

\screenshot{assets/images/Screenshot_2025-07-28_at_10.39.09_PM.png}{db\_nmap saving scan results to Metasploit database}

\begin{notebox}
\ic{setg} sets a global option that persists across all modules in the session.
\ic{set} only applies to the currently loaded module. Always use \ic{setg RHOSTS}
and \ic{setg RHOST} at the start to avoid repeating yourself.
\end{notebox}

% ─────────────────────────────────────────────────────────────────────────────
\section{Web Server Scanning}
\label{sec:web_scanning}

\begin{termbox}
# Find web directories and technologies
nmap --script http-enum -sV -p 80,443 target_ip

# Check for specific vulnerabilities
nmap -sV --script=http-enum -p 80 target_ip

# Run dirb for directory brute-force (if http-enum is slow)
dirb http://target_ip
dirb http://target_ip /usr/share/dirb/wordlists/common.txt
\end{termbox}

\screenshot{assets/images/Screenshot_2025-07-29_at_10.54.33_PM.png}{Web server enumeration using nmap http-enum script}

% ─────────────────────────────────────────────────────────────────────────────
\section{Learning Output}
\label{sec:scanning_learning}

\begin{takebox}
\textbf{What I Learned}
\begin{itemize}
  \item \ic{-sV -sC} together is my default scan. It gives version info and runs
        useful scripts automatically --- especially valuable for SMB, HTTP, and SSH.
  \item Running nmap from inside Metasploit saves time later. I don't have to
        manually re-enter IPs into modules --- the database already has everything.
  \item \ic{setg RHOST} and \ic{setg RHOSTS} at the start of a session is habit I
        built quickly after forgetting to set the IP on the third exploit attempt.
\end{itemize}

\textbf{Enumeration Mindset --- What to Check First}
\begin{enumerate}
  \item \ic{nmap -sn 192.168.x.0/24} --- confirm which hosts are alive
  \item \ic{nmap -sV -sC -p- target\_ip} --- full port scan with version and scripts
  \item Cross-reference open ports with known vulnerabilities in searchsploit
  \item Feed results into Metasploit workspace for tracking
\end{enumerate}

\textbf{Common Mistakes}
\begin{itemize}
  \item Scanning only the top 1000 ports and missing something on a high port
  \item Not running \ic{-sC} and missing SMB signing info or HTTP server details
        that NSE scripts provide automatically
  \item Forgetting to start PostgreSQL before msfconsole --- the DB won't connect
        and \ic{db\_nmap} won't save anything
\end{itemize}

\textbf{Real-World Impact}
\begin{itemize}
  \item Open port 445 with SMBv1 = immediate EternalBlue candidate
  \item Open port 80 with \ic{/webdav/} returning 401 = WebDAV attack surface
  \item Open port 3306 (MySQL) with no authentication = critical data exposure
\end{itemize}
\end{takebox}

\end{notebox}